\documentclass[10pt,a4paper]{amsart}
\usepackage[margin=19mm]{geometry}
\usepackage{amsmath,amssymb,mathtools}
\usepackage[scr=rsfs,cal=euler]{mathalfa}
\usepackage{xfrac}
\usepackage[unicode,hidelinks,bookmarks=false]{hyperref}
\newcommand{\ie}{i.e.\ }
\newcommand{\cf}{cf.\ }
\newcommand{\resp}{respectively\ }
\newcommand{\Subsection}{\subsection}
\providecommand{\nobreakdash}{\nobreak\hbox{-}}
\setlength{\emergencystretch}{2em}
\multlinegap0pt
\usepackage{bm}
\usepackage{bbm}
\usepackage{dsfont}
\usepackage{xspace}
\usepackage{cancel}
\usepackage{mathtools}
\usepackage{amsthm}
\makeatletter
\@ifundefined{lemma}{\newtheorem{lemma}{Lemma}}{}
\makeatother
\makeatletter
\expandafter\ifx\csname boxedalgo\endcsname\relax
\newenvironment{boxedalgo}
  {\begin{center}
  \begin{boxedminipage}{\linewidth}}
\fi
\providecommand{\cG}[0]{\ensuremath{\mathcal G}\xspace}
\providecommand{\cI}[0]{\ensuremath{\mathcal I}\xspace}
\providecommand{\conv}[1]{\ensuremath{\mathrm{conv}{\left({#1}\right)}\xspace}}
\providecommand{\convo}[1]{\ensuremath{\mathrm{conv}^o{\left({#1}\right)}\xspace}}
\newcommand{\fm}{\mathfrak{m}}
\newcommand{\fp}{\mathfrak{p}}
\newcommand{\fq}{\mathfrak{q}}
\newcommand{\fr}{\mathfrak{r}}
\let\geq\geqslant



%%% Local Commands
\providecommand{\defeq}[0]{\ensuremath{\;\coloneqq\;}\xspace}
\providecommand{\imm}[0]{\ensuremath{\mathrm{imm}}\xspace}
\let\leq\leqslant
\let\geq\geqslant



%%% Local Commands
\providecommand{\defeq}[0]{\ensuremath{\;\coloneqq\;}\xspace}
\makeatother
\title[Solving polynomial inequalities over spaces of convex sets a]{Solving polynomial inequalities over spaces of convex sets and applications}
\author{Saugata Basu, Hamidreza Amini Khorasgani, Hemanta K. Maji, Hai H. Nguyen}
\date{Source: arXiv:2608.07794v1 (CC BY 4.0)}
\begin{document}
\maketitle

\noindent\emph{Source and licence.} Solving polynomial inequalities over spaces of convex sets and applications. Version arXiv:2608.07794v1 (CC BY 4.0), distributed under the Creative Commons Attribution 4.0 International licence, as recorded in the arXiv licence metadata for this exact version and shown on the abstract page https://arxiv.org/abs/2608.07794v1. Full rights notice: licenses/arxiv-2608.07794-CC-BY-4.0.txt.\par\smallskip

\noindent\textit{Editorial scope.} Counted unit: the Lemma whose source label is \texttt{lem:root-immediacy-repair} in the article (source0.tex lines 4053--4064, source label \texttt{lem:root-immediacy-repair}), delivered together with the article's own complete original proof of it (source0.tex lines 4066--4464, one \texttt{proof} environment of 399 lines). No mathematics is written, repaired or re-derived: statement and proof are the article's own text line for line. Required context retained uncounted: lem:new (lines 3178--3230). Every definition, display and label of those lines is retained unchanged. Omitted from this package and not redistributed: 1--3177, 3231--4052, 4065--4065, 4465--5761. Nothing inside a retained excerpt is omitted or altered: every retained body is delivered exactly as the article's lines are written, so no omission receipt of any kind accompanies this package. Citations: the retained excerpts carry no citation command at all, so no bibliography entry of the article is transcribed and no reference is redistributed. Reference adaptation: none was needed and none was made; every reference inside the delivered unit resolves inside it, so no unresolved reference or citation is produced. Printed slips kept literal: no printed word, symbol, hypothesis or number of the counted statement or of its proof was altered. Layout and class: the article's own class is \texttt{amsart} with packages including inputenc, amsfonts, color, epsf, mathbbol, amsthm, amssymb, xspace, geometry, boxedminipage, enumerate, comment, todonotes, stmaryrd; the wrapper is the lane's pinned \texttt{amsart} editorial wrapper at 10pt on A4, which restates the theorem-like environments the retained text needs and adds the article's own macro definitions that the retained text needs (\texttt{\textbackslash cG}, \texttt{\textbackslash cI}, \texttt{\textbackslash conv}, \texttt{\textbackslash convo}, \texttt{\textbackslash fm}, \texttt{\textbackslash fp}, \texttt{\textbackslash fq}, \texttt{\textbackslash fr}, \texttt{\textbackslash geq}, \texttt{\textbackslash imm}, \texttt{\textbackslash leq}), with extra wrapper packages (bm, bbm, dsfont, xspace, cancel, mathtools). The delivered numbering is the unit's own. Assets: no retained span contains a figure environment, a graphics-inclusion command, a PDF-page inclusion, a table or a TikZ picture; no image file is copied into this package, the package declares no asset dependency and no image file is redistributed with it. Licence: version arXiv:2608.07794v1 of this article is distributed under the Creative Commons Attribution 4.0 International licence, as recorded in the arXiv licence metadata for this exact version and shown on the abstract page https://arxiv.org/abs/2608.07794v1. A full rights notice is delivered at licenses/arxiv-2608.07794-CC-BY-4.0.txt. Characters: every retained excerpt is pure ASCII, so the delivered engine, XeLaTeX with its default Latin Modern text font, reports no missing character.
\begin{lemma}
\label{lem:new}
Let \(v_1,\ldots,v_n \in 
V\), and let
\[
v=\sum_{i=1}^n \alpha_i v_i,
\qquad 
\alpha_i\geq 0,
\qquad
\sum_{i=1}^n \alpha_i=1.
\]
Let \(\mathcal S\) denote the collection of all subsets
\[
\sigma \subseteq \{1,\ldots,n\}
\]
such that the points \(\{v_i : i\in \sigma\}\) are affinely independent and
\[
v \in \operatorname{relint}
\operatorname{conv}\{v_i : i\in \sigma\}.
\]
For each \(\sigma\in\mathcal S\), let \(I_\sigma=(I_{\sigma,1},\ldots,I_{\sigma,n})\in\mathbb{R}^n\) be the coefficient vector defined by
\[
I_{\sigma,i}=0 \quad\text{if } i\notin \sigma,
\]
and, for \(i\in \sigma\), by the unique barycentric coordinates satisfying
\[
v=\sum_{i\in\sigma} I_{\sigma,i}v_i,
\qquad
\sum_{i\in\sigma} I_{\sigma,i}=1,
\qquad
I_{\sigma,i}>0.
\]
Then every convex representation of \(v\) as a convex combination of
\(v_1,\ldots,v_n\) is a convex combination of the vectors \(I_\sigma\).
In particular, there exist numbers \(t_\sigma\geq 0\), with
\[
\sum_{\sigma\in\mathcal S} t_\sigma=1,
\]
such that
\[
(\alpha_1,\ldots,\alpha_n)
=
\sum_{\sigma\in\mathcal S} t_\sigma I_\sigma.
\]
Equivalently,
\[
\alpha_i
=
\sum_{\sigma\in\mathcal S} t_\sigma I_{\sigma,i}
\qquad
\text{for every } i=1,\ldots,n.
\]
\end{lemma}

\begin{lemma}[Root-immediacy repair]
\label{lem:root-immediacy-repair}
Assume that $\dim W_i=1$ for all $1\leq i\leq k$.
Let $\Pi=(i_1,\ldots,i_\ell)$, and let $T$ be a
$(\Pi,\cG)$-gridded witness tree for a point $v$.
Assume that every node $n\in \mathrm{nodes}(T)$ with
\[
  0<d:=\mathrm{depth}_T(n)\leq \ell-1
\]
is $i_{d+1}$-immediate.  Then there exists a
$(\Pi,\cG)$-gridded immediate witness tree $\widetilde T$ for $v$.
\end{lemma}

\begin{proof}
We give the proof after relabelling the coordinates so that
\[
  \Pi=(1,\ldots,\ell).
\]
The general case is identical, replacing the coordinate $q$ below by
$i_q$ everywhere.

If $\ell=0$, then $T$ is already $(\Pi,\cG)$-gridded immediate, since
there are no prescribed levels at which immediacy must be checked.
Thus we may assume $\ell\geq 1$.

Let
\[
  \fr=\mathrm{root}(T), \qquad N=\mathrm{children}_T(\fr).
\]
Write
\[
  x_\fm=v(\fm),\qquad a_\fm=\alpha(\fm),\qquad \fm\in N.
\]
Since $T$ is a witness tree and $i(\fr)=1$, we have
\[
  v=\sum_{\fm\in N} a_\fm x_\fm,
\]
and for every $q\neq 1$ and every $\fm\in N$,
\[
  \pi_q(x_\fm)=\pi_q(v).
\]
Since $T$ is $(\Pi,\cG)$-gridded, every root child lies at
depth $1$, and therefore its first coordinate is grid-valued:
\[
  \pi_1(x_\fm)\in \cG^{(1)}\qquad (\fm\in N).
\]

Set
\[
  w_1:=\pi_1(v).
\]
Because $1\in \cI_{\cG}(v)$, we have $w_1\notin \cG^{(1)}$. Hence, since
$\dim W_1=1$, the immediate cell
\[
  \imm_1(w_1)
\]
has two elements. Write
\[
  \imm_1(w_1)=\{w_1^{(1)},w_1^{(2)}\}.
\]
There are unique positive numbers $\mu^{(1)},\mu^{(2)}$ such that
\[
  w_1=\mu^{(1)} w_1^{(1)}+\mu^{(2)}w_1^{(2)},\qquad
  \mu^{(1)}+\mu^{(2)}=1.
\]

We now decompose the old root weights conditionally over the two
immediate endpoints. Let
\[
  y_\fm:=\pi_1(x_\fm)\in \cG^{(1)}.
\]
Since
\[
  w_1=\sum_{\fm\in N} a_\fm y_\fm,
\]
we may apply Lemma~\ref{lem:new} to this convex representation.  Let
$\mathcal A$ be the collection of non-empty subsets
$\sigma\subseteq N$ such that the labeled set
$\{y_m:m\in\sigma\}$ is affinely independent and
\[
  w_1\in 
  \convo{\{y_\fm \;:\; \fm\in\sigma\}}.
\]
For $\fm\notin\sigma$, set $b_\sigma(\,\cdot\,,y_\fm)=0$.  Lemma~\ref{lem:new}  gives
numbers $t_\sigma\geq 0$, $\sigma\in\mathcal A$, satisfying
\[
  \sum_{\sigma\in\mathcal A}t_\sigma=1
\]
and, for each $\fm\in N$,
\begin{equation}\label{eq:old-weight-decomp}
  a_\fm
  =
  \sum_{\sigma\in\mathcal A}
  t_\sigma\, b_\sigma(w_1,y_\fm).
\end{equation}

For every $\sigma\in\mathcal A$, the points $w_1^{(1)}$ and
$w_1^{(2)}$ belong to
\[
  \conv{\{y_\fm \;:\; \fm\in\sigma\}}.
\]
Indeed, $\sigma$ is a one-dimensional simplex whose grid vertices
bracket $w_1$, and $w_1^{(1)},w_1^{(2)}$ are the two immediate grid
points surrounding $w_1$.

For $j=1,2$ and $m\in N$, define
\begin{equation}\label{eq:nu-def}
  \nu_\fm^{(j)}
  :=
  \sum_{\sigma\in\mathcal A}
  t_\sigma\, b_\sigma(w_1^{(j)},y_\fm).
\end{equation}
Then $\nu_\fm^{(j)}\geq 0$ and
\begin{equation}\label{eq:nu-sums}
  \sum_{\fm\in N}\nu_m^{(j)}=1,\qquad
  \sum_{\fm\in N}\nu_\fm^{(j)}y_\fm=w_1^{(j)}.
\end{equation}
Moreover, by affine linearity of barycentric coordinates and the
identity
\[
  w_1=\mu^{(1)}w_1^{(1)}+\mu^{(2)}w_1^{(2)},
\]
we have, for each $\fm\in N$,
\begin{align}
  a_\fm
  &=
  \sum_{\sigma\in\mathcal A}
  t_\sigma\,b_\sigma(w_1,y_\fm) \notag\\
  &=
  \sum_{\sigma\in\mathcal A}
  t_\sigma\,
  b_\sigma\bigl(\mu^{(1)}w_1^{(1)}+\mu^{(2)}w_1^{(2)},y_\fm\bigr)
  \notag\\
  &=
  \mu^{(1)}
  \sum_{\sigma\in\mathcal A}
  t_\sigma\,b_\sigma(w_1^{(1)},y_\fm)
  +
  \mu^{(2)}
  \sum_{\sigma\in\mathcal A}
  t_\sigma\,b_\sigma(w_1^{(2)},y_\fm) \notag\\
  &=
  \mu^{(1)}\nu_\fm^{(1)}+\mu^{(2)}\nu_\fm^{(2)}.
  \label{eq:a-nu}
\end{align}

We now construct the new tree $\widetilde T$.

First create a new root $\widetilde{\fr}$ with
\[
  v(\widetilde{\fr})=v,\qquad i(\widetilde{\fr})=1.
\]
Its two children are denoted $\widetilde{\fp}^{(1)}_{\emptyset}$ and
$\widetilde{\fp}^{(2)}_{\emptyset}$.  We set
\[
  \alpha(\widetilde{\fp}^{(j)}_{\emptyset})=\mu^{(j)}
\]
and
\begin{equation}\label{eq:root-child-value}
  v(\widetilde{\fp}^{(j)}_{\emptyset})
  :=
  \sum_{\fm\in N}\nu_\fm^{(j)}v(\fm),
  \qquad j=1,2.
\end{equation}
Then the new root convex identity holds:
\begin{align*}
  \sum_{j=1}^2 \alpha(\widetilde{\fp}^{(j)}_{\emptyset})
  v(\widetilde{\fp}^{(j)}_{\emptyset})
  &=
  \sum_{j=1}^2 \mu^{(j)}
  \sum_{\fm\in N}\nu_\fm^{(j)}v(\fm)  \\
  &=
  \sum_{\fm\in N}
  \bigl(\mu^{(1)}\nu_\fm^{(1)}+\mu^{(2)}\nu_\fm^{(2)}\bigr)v(\fm) \\
  &=
  \sum_{\fm\in N}a_\fm \cdot v(\fm)=v.
\end{align*}
Also, by \eqref{eq:nu-sums},
\[
  \pi_1(v(\widetilde{\fp}^{(j)}_{\emptyset}))=w_1^{(j)}.
\]
For $q\neq 1$, since all root children $m$ agree with $v$ in the
$W_q$-coordinate,
\[
  \pi_q(v(\widetilde{\fp}^{(j)}_{\emptyset}))
  =
  \sum_{\fm\in N}\nu_\fm^{(j)}\pi_q(v(\fm))
  =
  \pi_q(v).
\]
Thus the new root is a valid witness-tree split in direction $1$.
Furthermore, the map
\[
  \{\widetilde{\fp}^{(1)}_{\emptyset},\widetilde{\fp}^{(2)}_{\emptyset}\}
  \longrightarrow
  \imm_1(w_1),
  \qquad
  \widetilde{\fp}^{(j)}_{\emptyset}\longmapsto
  \pi_1(v(\widetilde{\fp}^{(j)}_{\emptyset}))
\]
is bijective. Hence the new root is $1$-immediate.

It remains to construct the subtrees below
$\widetilde{\fp}^{(1)}_{\emptyset}$ and $\widetilde{\fp}^{(2)}_{\emptyset}$.  This is done
synchronously by averaging corresponding nodes in the old subtrees
rooted at the old children $m\in N$.

We define, recursively, a family of words $\mathcal W_d$ for
$1\leq d\leq \ell$.  Let
\[
  \mathcal W_1:=\{\emptyset\},
\]
and for $\fm\in N$ set
\[
  \fm_{\emptyset}:=\fm.
\]
Suppose $\mathcal W_d$ has been defined for some $1\leq d<\ell$, and
let $\eta\in\mathcal W_d$.  The node $m_\eta$ has depth $d$ in the
old tree $T$.  Since $d\geq 1$ and $d\leq \ell-1$, the hypothesis
implies that $m_\eta$ is $(d+1)$-immediate.

We first observe the following invariant:

\[
  \text{for fixed }d,\eta,\text{ the points }v(\fm_\eta),\ \fm\in N,
  \text{ agree in every }W_q\text{-coordinate with }q\neq 1.
\]

For $d=1$, this is exactly the projection condition for the original
root split.  If the invariant holds for $d$, then all $\fm_\eta$ have
the same $W_{d+1}$-coordinate.  Hence the immediate cell
\[
  E_{d+1}(\eta):=
  \imm_{d+1}\bigl(\pi_{d+1}(v(\fm_\eta))\bigr)
\]
is independent of $\fm$.  Since $\fm_\eta$ is $(d+1)$-immediate, for each
$e\in E_{d+1}(\eta)$ there is a unique child of $\fm_\eta$, denoted
$m_{\eta,e}$, such that
\[
  \pi_{d+1}(v(\fm_{\eta,e}))=e.
\]
Define
\[
  \mathcal W_{d+1}
  :=
  \{(\eta,e):\eta\in\mathcal W_d,\ e\in E_{d+1}(\eta)\}.
\]
The invariant is preserved at level $d+1$: if $q=d+1$, then the
$W_q$-coordinate of all $v(\fm_{\eta,e})$ is the common endpoint $e$;
if $q\neq 1,d+1$, then the old witness-tree projection condition for
the split $\fm_\eta\to \fm_{\eta,e}$ gives
\[
  \pi_q(v(\fm_{\eta,e}))=\pi_q(v(\fm_\eta)),
\]
which is independent of $\fm$ by the induction hypothesis.

Now, for each $j=1,2$, each $1\leq d\leq \ell$, and each
$\eta\in\mathcal W_d$, introduce a new node $\widetilde{\fp}^{(j)}_\eta$ and define
\begin{equation}\label{eq:z-def}
  v(\widetilde{\fp}^{(j)}_\eta)
  :=
  \sum_{\fm\in N}\nu_\fm^{(j)}\cdot v(\fm_\eta).
\end{equation}
This agrees with \eqref{eq:root-child-value} when $d=1$ and
$\eta=\emptyset$.

If $d<\ell$, then $\widetilde{\fp}^{(j)}_\eta$ is declared to be an internal node
with
\[
  i(\widetilde{\fp}^{(j)}_\eta)=d+1.
\]
Its children are the nodes
\[
  \widetilde{\fp}^{(j)}_{\eta,e},
  \qquad e\in E_{d+1}(\eta).
\]
Let $\beta_{\eta,e}$ be the barycentric coordinate of the common point
$\pi_{d+1}(v(\fm_\eta))$ with respect to the immediate cell
$E_{d+1}(\eta)$.  These numbers are independent of $m$, and they are
precisely the weights of the old immediate split
\[
  \fm_\eta\longrightarrow \{\fm_{\eta,e}:e\in E_{d+1}(\eta)\}.
\]
Therefore
\[
  v(\fm_\eta)
  =
  \sum_{e\in E_{d+1}(\eta)}
  \beta_{\eta,e}v(\fm_{\eta,e})
\]
for every $\fm\in N$.  We set
\[
  \alpha(\widetilde{\fp}^{(j)}_{\eta,e})=\beta_{\eta,e}.
\]
Then
\begin{align*}
  \sum_{e\in E_{d+1}(\eta)}
  \alpha(\widetilde{\fp}^{(j)}_{\eta,e})v(\widetilde{\fp}^{(j)}_{\eta,e})
  &=
  \sum_{e\in E_{d+1}(\eta)}
  \beta_{\eta,e}
  \sum_{\fm\in N}\nu_\fm^{(j)}v(\fm_{\eta,e})       \\
  &=
  \sum_{\fm\in N}\nu_\fm^{(j)}
  \sum_{e\in E_{d+1}(\eta)}
  \beta_{\eta,e}v(\fm_{\eta,e})                 \\
  &=
  \sum_{\fm\in N}\nu_\fm^{(j)}v(\fm_\eta)            \\
  &=
  v(\widetilde{\fp}^{(j)}_\eta).
\end{align*}
Thus the convex identity holds at $\widetilde{\fp}^{(j)}_\eta$.

The projection condition also holds. Let $q\neq d+1$.  Then, for
each $m$ and each $e\in E_{d+1}(\eta)$,
\[
  \pi_q(v(\fm_{\eta,e}))=\pi_q(v(\fm_\eta)).
\]
Averaging over $\fm$ gives
\[
  \pi_q(v(\widetilde{\fp}^{(j)}_{\eta,e}))
  =
  \pi_q(v(\widetilde{\fp}^{(j)}_\eta)).
\]
So the split at $\widetilde{\fp}^{(j)}_\eta$ is a valid witness-tree split in
direction $d+1$.

Moreover,
\[
  \pi_{d+1}(v(\widetilde{\fp}^{(j)}_{\eta,e}))=e.
\]
Hence the map
\[
  \mathrm{children}(\widetilde{\fp}^{(j)}_\eta)\longrightarrow E_{d+1}(\eta),
  \qquad
  \widetilde{\fp}^{(j)}_{\eta,e}\longmapsto
  \pi_{d+1}(v(\widetilde{\fp}^{(j)}_{\eta,e}))
\]
is bijective.  Thus $\widetilde{\fp}^{(j)}_\eta$ is $(d+1)$-immediate.

After carrying out this construction for $d=1,\ldots,\ell-1$, we have
constructed all prescribed levels.  The nodes $\widetilde{\fp}^{(j)}_\eta$ with
$\eta\in\mathcal W_\ell$ lie at depth $\ell$.  They have all
$W$-coordinates in the grid:
coordinate $1$ is $w_1^{(j)}\in G^{(1)}$, coordinates
$2,\ldots,\ell$ are immediate endpoints, and all coordinates outside
$\{1,\ldots,\ell\}$ were already grid-valued in the original
$(\Pi,\cG)$-gridded tree.

We now attach the old lower subtrees. Fix $j\in\{1,2\}$ and
$\eta\in\mathcal W_\ell$.  Let
\[
  N_j:=\{\fm\in N:\nu_\fm^{(j)}>0\}.
\]
Declare $\widetilde{\fp}^{(j)}_\eta$ to be an internal node with index $1$, and give
it children $\widetilde{\fq}^{(j)}_{\eta,\fm}$ indexed by $\fm\in N_j$.  Set
\[
  \alpha(\widetilde{\fq}^{(j)}_{\eta,\fm})=\nu_\fm^{(j)},\qquad
  v(\widetilde{\fq}^{(j)}_{\eta,\fm})=v(\fm_\eta).
\]
Then
\[
  v(\widetilde{\fp}^{(j)}_\eta)
  =
  \sum_{\fm\in N_j}\nu_\fm^{(j)}v(\fm_\eta).
\]
For fixed $\eta$, the old nodes $\fm_\eta$, $\fm\in N$, agree in every
coordinate except possibly $W_1$.  Therefore this is a valid
witness-tree split in direction $1$.

Below each $\widetilde{\fq}^{(j)}_{\eta,\fm}$ we copy the old subtree of $T$ rooted
at $m_\eta$, with the only change that the incoming weight of the root
copy is now $\nu_\fm^{(j)}$. All outgoing weights, node values, node
indices, and descendant subtrees are copied from $T$. Leaves are
therefore still labeled by points of $S$.

This completes the construction of $\widetilde T$.

We verify that $\widetilde T$ has the required properties.

First, $\widetilde T$ is a witness tree. The root convex identity and
projection condition were checked above.  The same checks were carried
out for every synchronized node $\widetilde{\fp}^{(j)}_\eta$ at depths
$1,\ldots,\ell-1$.  The bottom attachment is a valid split in
coordinate $1$, and all lower copied subtrees inherit the witness-tree
identities from $T$.  The tree is finite because $T$ is finite and the
construction introduces only finitely many new averaged nodes.

Second, $\widetilde T$ is $(\Pi,\cG)$-gridded.  At depth $0$, this
holds because $\cI_\cG(v)=\{1,\ldots,\ell\}$.  At depth $d$ with
$1\leq d\leq \ell$, the construction has already fixed coordinates
$1,\ldots,d$ to grid endpoints, and all coordinates outside
$\{1,\ldots,\ell\}$ were grid-valued from the original gridded tree.
Thus every coordinate outside the remaining set
\[
  \{d+1,\ldots,\ell\}
\]
is grid-valued. For $d\leq \ell-1$, the node index is $d+1$ by
construction. At depths $d\geq \ell$, all coordinates are
grid-valued, either by the construction above or because the lower
subtree was copied from nodes of the original gridded tree at depth at
least $\ell$.

Third, $\widetilde T$ is immediate at all prescribed levels. The root
is $1$-immediate by construction.  For each $1\leq d\leq \ell-1$, the
children of every node $\widetilde{\fp}^{(j)}_\eta$ at depth $d$ are indexed
bijectively by the immediate cell in coordinate $d+1$.  Hence every
node at depth $d$ is $(d+1)$-immediate.

Therefore $\widetilde T$ is a $(\Pi,\cG)$-gridded immediate witness
tree for $v$.
\end{proof}

\end{document}
